% Part: first-order-logic
% Chapter: tableaux
% Section: provability-consistency

\documentclass[../../../include/open-logic-section]{subfiles}

\begin{document}

\iftag{FOL}
      {\olfileid{fol}{tab}{prv}}
      {\olfileid{pl}{tab}{prv}}

\olsection{\usetoken{S}{derivability} en consistentie: wat volgt en wat samen kan kloppen}

We bewijzen nu een paar eigenschappen van !!{derivability}: wanneer
een zin formeel uit aannamen volgt. Die eigenschappen zijn ook op
zichzelf interessant. Bovendien hebben we ze stuk voor stuk nodig in
het bewijs van de volledigheidsstelling.

\begin{prop}\ollabel{prop:provability-contr}
  Stel dat $\Gamma \Proves !A$ en dat $\Gamma \cup \{!A\}$
  inconsistent is. Dan is $\Gamma$ inconsistent.
\end{prop}

\begin{proof}
Er zijn twee eindige verzamelingen:
$\Gamma_0 = \{!B_1, \dots, !B_n\}$ en
$\Gamma_1 =\{!C_1, \dots, !C_n\} \subseteq \Gamma$. Daarvoor hebben
  \begin{align*}
    \{\sFmla{\False}{!A}, &
    \sFmla{\True}{!B_1}, \dots, \sFmla{\True}{!B_n}\} \\
    \{\sFmla{\True}{!A}, &
    \sFmla{\True}{!C_1}, \dots, \sFmla{\True}{!C_m}\}
  \end{align*}
  beide verzamelingen gesloten !!{tableau}s. Pas op $!A$ de regel
  \Cut{} toe. Daarmee kunnen we de twee tableaus samenvoegen tot één
  gesloten !!{tableau}. Dat bewijst dat
  $\Gamma_0 \cup \Gamma_1$ inconsistent is. Nu geldt
  $\Gamma_0 \subseteq \Gamma$ en $\Gamma_1 \subseteq \Gamma$. Daarom
  geldt ook $\Gamma_0 \cup \Gamma_1 \subseteq \Gamma$. Dus is
  $\Gamma$~inconsistent.
\end{proof}

\begin{prop}
\ollabel{prop:prov-incons}
$\Gamma \Proves !A$ geldt precies dan als
$\Gamma \cup \{\lnot !A\}$ inconsistent is.
\end{prop}

\begin{proof}
Neem eerst aan dat $\Gamma \Proves !A$. Dat betekent dat er een
gesloten !!{tableau} is voor
\[
\{\sFmla{\False}{!A},
\sFmla{\True}{!B_1}, \dots, \sFmla{\True}{!B_n}\}
\]
Pas nu de regel $\TRule{\True}{\lnot}$ toe. Zo maken we hiervan een
gesloten !!{tableau} voor
\[
\{\sFmla{\True}{\lnot !A},
\sFmla{\True}{!B_1}, \dots, \sFmla{\True}{!B_n}\}.
\]

Nu de andere richting. Neem een gesloten !!{tableau} voor de tweede
verzameling. Daar maken we als volgt een gesloten !!{tableau} voor de
eerste verzameling van. Verwijder elke formule die ontstaat door
\TRule{\True}{\lnot} toe te passen op de eerste aanname
$\sFmla{\True}{\lnot !A}$. Verwijder ook die aanname zelf. Voeg daarna
de aanname $\sFmla{\False}{!A}$ toe.
Waarom blijven alle takken gesloten? Een tak kan al gesloten zijn
doordat hij de conclusie bevat van \TRule{\True}{\lnot}, toegepast op
$\sFmla{\True}{\lnot !A}$. Die conclusie is
$\sFmla{\False}{!A}$. De bijbehorende tak in het nieuwe !!{tableau} is
dus ook gesloten. Een tak in het oude tableau kan ook gesloten zijn
doordat hij zowel de aanname $\sFmla{\True}{\lnot !A}$ als
$\sFmla{\False}{\lnot !A}$ bevat. Pas dan
$\TRule{\False}{\lnot}$ toe op $\sFmla{\False}{\lnot !A}$. Dat geeft
$\sFmla{\True}{!A}$. Ook deze tak sluit, want we hebben
$\sFmla{\False}{!A}$ als aanname toegevoegd.
\end{proof}

\begin{prob}
Geef een bewijs dat $\Gamma \Proves \lnot !A$ precies dan geldt als
$\Gamma \cup \{!A\}$ inconsistent is.
\end{prob}

\begin{prop}\ollabel{prop:explicit-inc}
  Stel dat $\Gamma \Proves !A$ en $\lnot !A \in \Gamma$. Dan is
  $\Gamma$ inconsistent.
\end{prop}

\begin{proof}
  Neem aan dat $\Gamma \Proves !A$ en $\lnot !A \in \Gamma$. Dan zijn
  er $!B_1$, \dots, $!B_n \in \Gamma$ waarvoor \
  \[
  \{\sFmla{\False}{!A}, \sFmla{\True}{!B_1}, \dots, \sFmla{\True}{!B_n}\}
  \]
  een gesloten tableau bestaat. Vervang de aanname
  \sFmla{\False}{!A} door \sFmla{\True}{\lnot !A}. Zet na de aannamen
  ook de conclusie van \TRule{\True}{\lnot}, toegepast op
  \sFmla{\False}{!A}. Verwees een !!{sentence} in het oude
  !!{tableau} voor haar rechtvaardiging naar regel~$1$? Dan verwijst
  die zin in het nieuwe !!{tableau} naar regel~$n+1$. Was het oude
  !!{tableau} gesloten, dan is het nieuwe dus ook gesloten. Alle
  aannamen liggen in~$\Gamma$. Het nieuwe tableau laat daarom zien dat
  $\Gamma$ inconsistent is.
\end{proof}

\begin{prop}\ollabel{prop:provability-exhaustive}
  Stel dat $\Gamma \cup \{!A\}$ en $\Gamma \cup \{\lnot !A\}$ beide
  inconsistent zijn. Dan is $\Gamma$ inconsistent.
\end{prop}

\begin{proof}
  Stel dat er $!B_1$, \dots, $!B_n \in \Gamma$ en $!C_1$, \dots,
  $!C_m \in \Gamma$ zijn waarvoor
  \begin{align*}
    \{\sFmla{\True}{!A}, &
    \sFmla{\True}{!B_1}, \dots, \sFmla{\True}{!B_n}\} \text{ and}\\
    \{\sFmla{\True}{\lnot !A}, &
    \sFmla{\True}{!C_1}, \dots, \sFmla{\True}{!C_m}\}
  \end{align*}
  allebei een gesloten !!{tableau} hebben. We bouwen daar één
  gecombineerd !!{tableau} van dat laat zien dat $\Gamma$ inconsistent
  is. Neem als aannamen
  $\sFmla{\True}{!B_1}$, \dots, $\sFmla{\True}{!B_n}$ samen met
  $\sFmla{\True}{!C_1}$, \dots, $\sFmla{\True}{!C_m}$. Pas daarna de
  regel \Cut{} toe. Er ontstaan twee takken. De ene begint met
  $\sFmla{\True}{!A}$ en de andere met $\sFmla{\False}{!A}$.  
  
  Voeg links het deel van het eerste !!{tableau} onder de aannamen
  toe. Alle regels zijn daar nog steeds juist toegepast. Elke aanname
  van het eerste !!{tableau} is immers beschikbaar, ook
  $\sFmla{\True}{!A}$. Daarom sluit elke tak onder
  $\sFmla{\True}{!A}$. 
  
  Voeg rechts het deel van het tweede !!{tableau} onder zijn aanname
  toe. Laat daar alle uitkomsten weg van toepassingen van
  $\TRule{\True}{\lnot}$ op $\sFmla{\True}{\lnot !A}$. De conclusie
  van $\TRule{\True}{\lnot}$ op $\sFmla{\True}{\lnot !A}$ is
  $\sFmla{\False}{!A}$. Die formule is toch al beschikbaar: ze staat
  rechts als conclusie van de regel \Cut{} in het gecombineerde
  !!{tableau}.

  Er blijven twee mogelijkheden over. Misschien sloot een tak in het
  tweede tableau doordat hij zowel de aanname
  $\sFmla{\True}{\lnot !A}$ bevatte als
  $\sFmla{\False}{\lnot !A}$. De eerste staat niet meer als aanname in
  het gecombineerde !!{tableau}. Pas daarom
  $\TRule{\False}{\lnot}$ toe op $\sFmla{\False}{\lnot !A}$. Dat geeft
  $\sFmla{\True}{!A}$. Dezelfde tak in het gecombineerde !!{tableau}
  sluit nu ook, want rechts staat al de conclusie van de regel \Cut{},
  namelijk $\sFmla{\False}{!A}$. Sloot een tak in het tweede
  !!{tableau} om een andere reden? Dan sluit de bijbehorende tak in het
  gecombineerde !!{tableau} ook. Elke !!{signed formula} behalve
  $\sFmla{\True}{\lnot !A}$ die op de oude tak stond, staat namelijk
  ook op de bijbehorende tak in het gecombineerde !!{tableau}.
  \end{proof}

\end{document}
